\subsection{An exact polynomial certificate for a new vanishing degree}

The remaining finite computation can be expressed without choosing
algebraic eigenvalues individually. Retain the entire indexed trace
multiset from the proof of VK-EXP7, including repeated values.

\begin{proposition}[VK-POLY-CERT]
Let $e=2k>0$ and let $p\ge7$ be a regular auxiliary prime. Put
$q=\operatorname{Nm}(w)$ for the chosen place of
$\mathbb Q(\sqrt{-3})$ above $p$, and define
\[
 F_{p,e}(X)=\prod_{s\in S_{p,e}}(X-q^e t_s),\qquad
 E_{p,e}=q^{2e}+q^e+1.
\]
Then $F_{p,e}$ is monic in $\mathbb Z[X]$, its value at $E_{p,e}$
is nonzero, and
\[
 |K_{2e}(\mathbb Z)|\mid |F_{p,e}(E_{p,e})|^2.
 \tag{K18a}
\]
More generally, let $P_{p,e}\in\mathbb Z[X]$ be monic and nonzero,
with $F_{p,e}\mid P_{p,e}$ in $\mathbb Q[X]$. Remove the maximal
power of $X-E_{p,e}$ from $P_{p,e}$, obtaining $P^*_{p,e}$.
Then
\[
 P^*_{p,e}(E_{p,e})\ne0,\qquad
 |K_{2e}(\mathbb Z)|\mid |P^*_{p,e}(E_{p,e})|^2.
 \tag{K18b}
\]
For two such auxiliary primes the order divides the square of the
gcd of these two nonzero evaluations. If every prime factor of this
gcd is below $2^{31}$, the already verified Vandiver range implies
$K_{2e}(\mathbb Z)=0$.
\end{proposition}

\begin{proof}
The source properties \texttt{TrData.int} and
\texttt{TrData.conj} in \path{AutomorphicVarieties/AutomorphicForms/Classifies.lean} say that every
$q^e t_s$ is an algebraic integer and that every Galois automorphism
permutes the indexed multiset. The elementary symmetric
coefficients are therefore rational algebraic integers. This gives
the asserted integral monic polynomial. Its evaluation satisfies
\[
 n_p(e)=(-1)^{|S_{p,e}|}F_{p,e}(E_{p,e}).
\]
The bound $|t_s|\le3$ and $q\ge7$ imply
$E_{p,e}>3q^e$, so every evaluated factor is nonzero.
The all-prime divisibility in VK-EXP7, using regularity to remove
the auxiliary prime itself, gives (K18a).

Monic polynomial division gives $P_{p,e}=F_{p,e}Q$ with
$Q\in\mathbb Z[X]$. Each exact division by the monic linear
polynomial $X-E_{p,e}$ also stays in $\mathbb Z[X]$.
Since $F_{p,e}(E_{p,e})\ne0$, $F_{p,e}$ has no such linear
factor. Unique factorization in $\mathbb Q[X]$ shows that it still
divides $P^*_{p,e}$; monic division makes the quotient integral.
This proves (K18b). Taking valuations at each prime proves the
gcd assertion. The final assertion uses the already cited vanishing
of every primary part at primes below $2^{31}$.
\end{proof}

Thus a proved, multiplicity-preserving inclusion of the source
trace system in a larger integral Hecke model is sufficient.
Removing the exact scalar root above can avoid a separate
Eisenstein projection. Taking the squarefree part is not permitted:
distinct representations may have the same trace at one prime,
and their multiplicities occur in the source integer.

Both $7$ and $13$ split in $\mathbb Q(\sqrt{-3})$ and are regular.
The following exact evaluation points were checked by integer
arithmetic; no corresponding Hecke polynomials were computed.
\[
\begin{array}{c|r|r|r}
K\text{-group}&e&p&E_{p,e}\\\hline
K_{12}&6&7&13\,841\,404\,851\\
K_{12}&6&13&23\,298\,089\,949\,291\\
K_{16}&8&7&33\,232\,936\,334\,403\\
K_{16}&8&13&665\,416\,609\,998\,910\,563
\end{array}
\]
The required highest weights are $(5,0,-5)$ and $(7,0,-7)$.
The open input is a complete integral Hecke model at the source's
specified levels, together with the coverage and normalization
proof. A dimension count or a spectrum computed from a
conjectural trace formula does not supply that input. The
literature and available-software checks, and this precise finite
certificate reduction, are recorded in
\path{evidence/work/k_numeric_followup.md}.
